\section{Score 与 Fisher 信息衡量局部可辨认性}
\label{sec:11-Mathematical-Statistics/03-Sufficient-Statistics-and-Information/02}

一次优惠券实验：五万次曝光，记录每次的折扣力度 \(d_i\) 与是否下单，拟合 logistic 回归

\[
p_i=\frac{1}{1+e^{-(\beta_0+\beta_1 d_i)}},
\]

想同时得到基线购买倾向 \(\beta_0\) 与折扣敏感度 \(\beta_1\)。结果是优化器跑了两次给出两组差很远的系数，标准误大到区间横跨零点，而两个系数的相关系数是 \(-0.998\)。工程侧先怀疑求解器，换了优化算法，换了初值，症状照旧。

翻回数据才看到问题：这五万次曝光发的是同一档折扣。所有 \(d_i\) 相等时，\(\beta_0+\beta_1 d\) 这一个数可以由无穷多组 \((\beta_0,\beta_1)\) 拼出来，高基线配低敏感与低基线配高敏感对数据的解释一模一样。似然沿着某个方向是平的，优化器落在哪儿全看初值。

这类故障不能靠看点估计发现，得看似然函数在真值附近的形状。

\begin{center}
\begin{tikzpicture}[x=1.2cm,y=0.78cm,font=\small,>=stealth]
  \draw[->,black!55] (-3.2,-3.2) -- (3.9,-3.2) node[right,font=\footnotesize] {\(\theta\)};
  \draw[dashed,black!60] (0,-4.7) -- (0,0.3);
  \node[font=\footnotesize,black!75,anchor=south] at (0,0.3) {真值 \(\theta_0\)};
  \draw[dashed,black!35] (-3.1,-1) -- (3.1,-1);
  \node[font=\footnotesize,black!60,anchor=south east] at (3.15,-0.95) {\(\ell\) 下降同样多};
  \draw[very thick,blue!55!black,domain=-1.15:1.15,samples=61,smooth] plot (\x,{-2.2*\x*\x});
  \draw[blue!55!black,domain=-2.9:2.9,samples=81,smooth] plot (\x,{-0.35*\x*\x});
  \node[font=\footnotesize,anchor=west] at (1.2,-2.9) {\(\ell\) 尖};
  \node[font=\footnotesize,anchor=west] at (2.98,-2.94) {\(\ell\) 平};
  \draw[dotted,black!55] (-0.674,-1) -- (-0.674,-3.6);
  \draw[dotted,black!55] (0.674,-1) -- (0.674,-3.6);
  \draw[dotted,black!55] (-1.69,-1) -- (-1.69,-4.5);
  \draw[dotted,black!55] (1.69,-1) -- (1.69,-4.5);
  \draw[<->,very thick,blue!55!black] (-0.674,-3.6) -- (0.674,-3.6);
  \node[font=\footnotesize,anchor=west] at (0.78,-3.6) {参数被按在窄区间里};
  \draw[<->,blue!55!black] (-1.69,-4.5) -- (1.69,-4.5);
  \node[font=\footnotesize,anchor=west] at (1.78,-4.5) {同样的数据说不清};
\end{tikzpicture}

\emph{同一个真值、同样的对数似然降幅，尖的那条把参数锁进窄区间，平的那条留下一大片说得通的取值。}
\end{center}

图里两条曲线的峰在同一个位置，也就是说两种情形下最优点估计可以完全一样。差别只在峰的两侧：粗线陡，参数稍微挪一点，数据的解释力马上变差；细线缓，挪出去很远似然也没掉多少。优惠券实验的病灶就是细线，只不过发生在二维参数空间的某一个方向上。Fisher 信息要做的，是把``峰有多尖''变成一个可以计算、可以在实验开跑前估出来的数。

\subsection{Score 记录参数微调时似然的反应}

对单个观测 \(X\sim p_\theta\)，score 函数是对数似然对参数的导数：

\[
U_\theta(X)=\frac{\partial}{\partial\theta}\log p_\theta(X).
\]

记号上留个心：上一节里 \(U(X)\) 指任意一个估计量，从这里起 \(U_\theta\) 专指 score。

它读的是一个局部量。某次观测算出的 \(U_\theta\) 绝对值大，说明在这条数据看来，\(\theta\) 挪一点点，模型对它的解释就会明显变好或明显变差；绝对值接近零，说明这条数据对 \(\theta\) 往哪边调没有意见。单条 score 本身不是可用的指标，因为它随数据变；能用的是它在整个分布下的表现。

\subsection{Score 在真值处平均为零}

在可以交换求导与积分、支持集不随参数变化的正则条件下，

\[
\begin{aligned}
\E_\theta[U_\theta(X)]
&=\int\frac{\partial}{\partial\theta}\log p_\theta(x)\cdot p_\theta(x)\,dx\\
&=\int\frac{\partial}{\partial\theta}p_\theta(x)\,dx
=\frac{\partial}{\partial\theta}\int p_\theta(x)\,dx
=\frac{\partial}{\partial\theta}1=0.
\end{aligned}
\]

骨架只有三步：把对数导数还原成密度导数，把导数挪到积分外，对常数 \(1\) 求导。条件对一遍——第二步要求交换求导与积分合法，第三步要求积分区域不随 \(\theta\) 移动。Uniform\((0,\theta)\) 这种支持集会动的模型两条都不满足，本节所有结论对它都不成立。

期望为零这件事说的是：站在真值上，数据对参数的平均推力恰好抵消。有些观测把估计往上推，有些往下拉，长期看没有系统性方向。于是可以拿来度量的就只剩波动幅度——推力越是忽大忽小，说明单条数据对参数的位置越敏感。

\subsection{Fisher 信息既是 score 的方差，也是似然的平均曲率}

Fisher 信息定义为 score 的二阶矩，由上一小节它同时就是方差：

\[
I(\theta)=\E_\theta\bigl[U_\theta(X)^2\bigr]=\Var_\theta\bigl(U_\theta(X)\bigr).
\]

向量参数下推广成矩阵 \(I(\theta)=\E_\theta[U_\theta U_\theta^{\trans}]\)，第 \((j,k)\) 个元素是两个分量 score 乘积的期望。

这个定义读起来还是关于斜率的，跟图里的``尖''对不上。把两者接起来只需要再求一次导。对恒等式 \(\E_\theta[U_\theta]=0\) 关于 \(\theta\) 求导，仍假定可以交换，并用 \(\partial_\theta p_\theta=U_\theta\,p_\theta\)：

\[
0=\int\Bigl[\frac{\partial^2}{\partial\theta^2}\log p_\theta(x)+\Bigl(\frac{\partial}{\partial\theta}\log p_\theta(x)\Bigr)^2\Bigr]p_\theta(x)\,dx,
\]

移项即得

\[
I(\theta)=-\E_\theta\Bigl[\frac{\partial^2}{\partial\theta^2}\log p_\theta(X)\Bigr].
\]

右边是对数似然二阶导的负期望，也就是平均曲率。图里粗线的曲率大，对应的 \(I(\theta)\) 就大。两种写法在正则模型下等价，遇到边界参数、参数相关的支持集或非光滑似然时得逐条重查条件，不能随手互换；模型设定错误时它们也会分家，后面还会回到这一点。

\subsection{一条 Bernoulli 样本值多少信息，取决于 \(\theta\) 落在哪}

单个 Bernoulli 观测的对数似然是 \(X\log\theta+(1-X)\log(1-\theta)\)，score

\[
U_\theta(X)=\frac{X}{\theta}-\frac{1-X}{1-\theta}=\frac{X-\theta}{\theta(1-\theta)}.
\]

分子的方差是 \(\theta(1-\theta)\)，平方取期望后

\[
I_1(\theta)=\frac{\theta(1-\theta)}{\theta^2(1-\theta)^2}=\frac{1}{\theta(1-\theta)}.
\]

独立样本的 score 是各条 score 之和，交叉项协方差为零，因此 \(I_n(\theta)=n\,I_1(\theta)\)。信息随样本量线性增长，不确定范围按 \(1/\sqrt n\) 收窄，这和经验一致。

值得盯一眼的是那个因子。\(\theta=0.5\) 处信息最小，一条样本给 \(4\) 个单位；\(\theta\) 挪向 \(0.01\)，一条样本给 \(101\) 个单位。同一份日志、同一个样本量，估一个接近 \(0.5\) 的转化率和估一个接近 \(0.01\) 的转化率，似然曲面的尖锐程度差了二十多倍——不过这里的信息是关于概率尺度上的绝对误差说的，若关心的是相对误差，罕见事件反而更难估，稍后重参数化那一小节会说明这两种读法为何不矛盾。至于 \(\theta\) 真的取到 \(0\) 或 \(1\)，公式给出的无穷大不能当真，那是非规则边界，正态渐近与下一节的下界都不适用。

\subsection{矩阵的秩暴露了哪个方向问不清}

回到优惠券实验。第 \(i\) 次曝光对信息矩阵的贡献是

\[
I_i(\beta)=p_i(1-p_i)\begin{pmatrix}1&d_i\\ d_i&d_i^2\end{pmatrix},
\]

前面的标量是 Bernoulli 那一项在 \(p_i\) 处的信息经链式法则折算的结果，后面的矩阵由这次曝光的输入方向 \((1,d_i)^{\trans}\) 决定。全部 \(d_i\) 相等时，每一项都是同一个秩为一的矩阵的倍数，加起来仍然秩一。信息矩阵奇异，沿着某个方向一阶敏感度为零，标准误在那个方向上没有定义——五万条记录改成五百万条也不会变。开头那次``求解器有问题''的排查，方向从一开始就错了。

同一个公式还给出另外两条判断。若折扣水平选在几乎人人下单或几乎无人下单的区段，\(p_i(1-p_i)\) 趋近零，即便方向矩阵满秩，每条记录的贡献也接近零，这是花了预算却没买到信息的典型方式。若把折扣换成药物剂量，数学结构一字不变，但可选的剂量范围受安全约束限死，Fisher 信息此时是比较备选方案的工具，它能告诉你哪种设计的可辨认性更高，不能替你取消约束。信息量取决于样本量，也取决于输入覆盖了参数空间的哪些方向，而后者往往是实验开跑前就已经定死的。

\subsection{观测信息是这一份数据的曲率，Fisher 信息是它的平均}

真正拿到手的只有一份数据 \(x\)。它的观测信息是对数似然在该点的负二阶导

\[
J(\theta;x)=-\frac{\partial^2}{\partial\theta^2}\ell(\theta;x),
\]

不含任何期望运算，直接由这份数据算出。上面那条恒等式说的正是 \(I(\theta)=\E_\theta[J(\theta;X)]\)：Fisher 信息是观测信息在所有可能数据上的平均。图里的两条曲线因此有两种读法，读成手上这一次实际得到的似然曲线，量的是 \(J\)；读成同一机制下重复实验的平均形状，量的是 \(I\)。

极大似然估计的协方差常用 \(J(\hat\theta;x)^{-1}\) 或 \(I(\hat\theta)^{-1}\) 估计，大样本下接近，有限样本下可以差出不少。\(J\) 如实反映手上这份数据碰巧有多陡，\(I\) 做了平均因而更稳，代价是它假定模型没设定错。模型确实错了的时候，score 的方差与负 Hessian 的期望不再相等，正确的协方差要写成三明治形式 \(H^{-1}JH^{-1}\)，其中 \(H\) 取负 Hessian 的期望或观测值，\(J=\E[U_\theta U_\theta^{\trans}]\)。模型正确时 \(H=J=I\)，三明治退回 \(I^{-1}\)；模型错误时两层各自有独立含义，这也是稳健标准误的来源。

\subsection{换一把尺子量参数，数值会变，几何不变}

令 \(\eta=g(\theta)\) 换一种参数化，链式法则给出 \(U_\eta=U_\theta\,d\theta/d\eta\)，平方取期望后

\[
I_\eta(\eta)=I_\theta(\theta)\Bigl(\frac{d\theta}{d\eta}\Bigr)^2.
\]

用概率、用 logit \(\eta=\log\bigl(\theta/(1-\theta)\bigr)\)、还是用百分比来表达同一个参数，Fisher 信息的数值不同。上一小节里罕见事件``信息更多''的说法，正是绑在概率这把尺子上的；换成 logit 尺度，\(\theta\) 趋近 \(0\) 时信息反而趋近零，与``罕见事件难估''的经验对上。所以引用一个 Fisher 信息的数值，必须同时交代参数用什么刻度写。

数值会变，但组合量 \(I(\theta)\,(d\theta)^2\) 在重参数化下不变，它像一把嵌在参数空间里的局部尺子。这一点在 KL 散度的二阶展开中兑现：相邻两个参数值满足

\[
\KL\bigl(p_\theta\,\Vert\,p_{\theta+d\theta}\bigr)\approx\tfrac12\,I(\theta)\,(d\theta)^2,
\]

于是 Fisher 信息量的是参数的微小移动在分布空间里走出了多远。信息大，两个靠得很近的参数值在数据分布上分得开，有限样本就能区分；信息小，它们在数据看来几乎是同一个分布，图里的细线画的就是这种情形。把这把尺子当成优化时的度量，梯度下降就变成自然梯度，详见\hyperref[sec:12-Real-Analysis-Support/09-Information-Geometry/02]{``自然梯度沿分布几何下降''}；几何一侧的完整说法见\hyperref[sec:12-Real-Analysis-Support/09-Information-Geometry/03]{``KL 散度与 Fisher 度量的局部关系''}。

矩阵版本的退化也有了几何解释。\(I(\theta)\) 有零特征值，意味着沿对应方向移动参数在一阶上不改变分布，模型在该方向不可识别。成因不止一种：参数冗余、标签交换对称、或者像优惠券实验那样输入没有覆盖到相应方向。检查特征值与条件数能发现弱识别，但别指望某个阈值一刀切——优化里的 Hessian 病态还掺着样本噪声与参数化的影响，与统计上的不可识别有关联而不等同。

图里那条粗线还剩一句话没兑现。峰越尖，参数被按得越紧，这句话到目前为止仍是定性的：\(I(\theta)\) 大，不确定范围小，但小到多少？下一节把它写成一个不等式，并且说明为什么那是一条谁也绕不过去的下界。
